\section*{\texorpdfstring{\S\CurrentExerciseGroup}{第{\CurrentExerciseGroup}节}}
\phantomsection
\addcontentsline{toc}{section}{第{\CurrentExerciseGroup}节习题}

\begin{enumerate}
\def\labelenumi{\arabic{enumi})}
\tightlist
\item
  证明：若 \(f\) 和 \(g\) 是两个数值函数 \(\geqslant 0\) ，则
\end{enumerate}

\[
\mu^ {\bullet} (\sup (f, g)) + \mu^ {\bullet} (\inf (f, g)) \leqslant \mu^ {\bullet} (f) + \mu^ {\bullet} (g)
\]

（参见 Ch. IV, §1, 习题 1）。

\begin{enumerate}
\def\labelenumi{\arabic{enumi})}
\setcounter{enumi}{1}
\tightlist
\item
  证明：对每个函数 \(f \geqslant 0\) （定义在 \(T\) 上），映射 \(\mu \mapsto \mu^{\bullet}(f)\) （\(\mathcal{M}_{+}(T)\) 到 \(\overline{\mathbf{R}}\) 中的映射）关于拟强拓扑是下半连续的（Ch. III, §1, 习题 8；参见 Ch. IV, §1, 习题 6 b)）。在什么条件下此映射连续？
\end{enumerate}

¶ 3) a) 设 \((f_{\alpha})_{\alpha \in \mathbb{A}}\) 是数值函数 \(\geqslant 0\) 组成的一个族，按关系 \(\leqslant\) 有向，且使映射 \(t \mapsto (f_{\alpha}(t))\) （\(\mathbf{T}\) 到 \(\mathbf{R}^{\mathbf{A}}\) 中的映射）是 \(\mu\) -可测的。设 \(f\) 是族 \((f_{\alpha})\) 的上包络。则 \(f\) 是 \(\mu\) -可测的，且

\[
\int^ {\bullet} f d \mu = \sup _ {\alpha \in A} \int^ {\bullet} f _ {\alpha} d \mu .\tag{1}
\]

\begin{enumerate}
\def\labelenumi{\alph{enumi})}
\setcounter{enumi}{1}
\tightlist
\item
  设 \((f_{\alpha})_{\alpha\in\mathrm{A}}\) 是定义在 T 上的数值函数 \(\geqslant0\) 组成的一个族，且具有下列性质：对 T 的每个紧子集 K 和每个 \(\varepsilon>0\) ，存在紧集 \(K'\subset K\) ，使得 \(\mu(K-K')\leqslant\varepsilon\) ，且在 \(K'\) 上所有函数 \(f_{\alpha}\) 的限制都是下半连续的。设 f 是族 \((f_{\alpha})\) 的上包络。证明 f 可测并满足关系式 (1)。
\end{enumerate}

给出一个映射 \(t \mapsto (f_{\alpha}(t))\) 的例子，它满足上述条件但不是 \(\mu\) -可测的。

\begin{enumerate}
\def\labelenumi{\arabic{enumi})}
\setcounter{enumi}{3}
\tightlist
\item
  设 \(\mathbf{T}\) 和 \(\mu\) 是 Ch. IV, §1 习题 5 中定义的局部紧空间和测度。
\end{enumerate}

\begin{enumerate}
\def\labelenumi{\alph{enumi})}
\item
  证明 \(\mu\) 是具有限支集测度序列的上确界，但 \(\mu^{*} \neq \mu^{\bullet}\) ，且 \(\mu\) 不是适度的。由此推出，当把其中的符号 \(\int^{\bullet}\) 换成 \(\int^{*}\) 时，即使集 A 可数，命题 11 也变为不真。
\item
  对每个实数 \(y\) ，用 \(f_y\) 表示化为点 \((0, y)\) （\(\mathbf{T}\) 的点）的单点集的特征函数。证明映射 \(t \mapsto (f_y(t))_{y \in \mathbf{R}}\) 是 \(\mu\) -可测的，但若 \(f = \sup_{y \in \mathbf{R}} f_y\) ，则 \(\mu^*(f) = +\infty\) 而 \(\sup_{y \in \mathbf{R}} \mu^*(f_y) = 0\) 。
\end{enumerate}

\begin{enumerate}
\def\labelenumi{\arabic{enumi})}
\setcounter{enumi}{4}
\tightlist
\item
  设 \((\mu_n)\) 是 \(T\) 上的正测度序列，拟强收敛于 \(\mu\) （在 \(\mathcal{M}(T)\) 中）（Ch. III, §1, 习题 8）。
\end{enumerate}

\begin{enumerate}
\def\labelenumi{\alph{enumi})}
\item
  证明：若映射 \(f\) （\(\mathbf{T}\) 到一个拓扑空间 \(\mathbf{G}\) 中的映射）是 \(\mu_n\) -可测的（对每个 \(n\) ），则它是 \(\mu\) -可测的（利用 Ch. IV, §1 的习题 6 b)）。
\item
  设 \(\mathbf{f}\) 是 T 到一个 Banach 空间 \(F\) 中的映射，它是本质 \(\mu_n\) -可积的（对每个 \(n\) ）。证明：若序列 \((\mu_n^\bullet (|\mathbf{f}|))\) 有界，则 \(\mathbf{f}\) 是本质 \(\mu\) -可积的（习题 2）。用一个例子说明不一定有 \(\int \mathbf{f} d\mu = \lim_{n\to \infty}\int \mathbf{f} d\mu_n\) （对 \(\mu_n\) 取紧空间上的一个点测度，使得 \(\lim \| \mu_n\| = 0\) ）；但当 \(\mathbf{f}\) 有界且具紧支集时，这一关系式成立。
\end{enumerate}

\begin{enumerate}
\def\labelenumi{\arabic{enumi})}
\setcounter{enumi}{5}
\tightlist
\item
  设 \(\mathfrak{K}\) 是 T 的紧子集组成的一个 \(\mu\) -稠密族。设 \(\mathbf{f}\) 是定义在 T 上、取值于一个 Banach 空间 F 中的函数，使得 \(\mathbf{f}\varphi_{\mathrm{K}}\) 是 \(\mu\) -可积的（对每个 \(\mathrm{K} \in \mathfrak{K}\) ）。若积分 \(\int \mathbf{f}\varphi_{\mathrm{K}} d\mu\) 沿有向集 \(\mathfrak{K}\) （按 \(\subset\) ）在 F 中有一个极限，则 \(\mathbf{f}\) 是本质 \(\mu\) -可积的。
\end{enumerate}

\setcounter{exercisegroup}{2}
\edef\CurrentExerciseGroup{\arabic{exercisegroup}}
